% !PRINTED 264
\weberpage{264}
Hiernach lassen sich leicht die Grade der Elemente unserer
Gruppe bestimmen. Wir haben dabei drei Fälle zu unterscheiden.

I.~\gesperrt{$a + d = \pm 2$.} \gesperrt{Jede Substitution dieser Art lässt sich
in die Normalform $T$ transformiren. Es ist aber für jeden
Exponenten $r$}
\[
T^r = \begin{pmatrix} 1, & r\,t \\ 0, & 1 \end{pmatrix},
\]
\gesperrt{und folglich ist der Grad dieser Substitution $p$, ausser
wenn $t = 0$, also $T$ die identische Substitution ist.}

Die Anzahl der Elemente dieser Art ist leicht zu bestimmen.
Nehmen wir zunächst $c = 0$, so können wir $b$ beliebig wählen
und für $a, d$ erhalten wir aus $a + d = \pm 2$, $a\,d = 1$ die beiden
Bestimmungen $a = d = \pm 1$. Dies giebt $2p$ Formen, die identische
 Substitution eingeschlossen. Nehmen wir sodann $a$ und $c$
beliebig, jedoch $c$ von Null verschieden, was $p\,(p-1)$ Möglichkeiten
 giebt, so folgt aus $a\,d - b\,c = 1$
\[
b = -c^{-1} + a\,d\,c^{-1},
\]
und zu jedem $a$ kann $d$ auf zwei Arten bestimmt werden. Die
Gesammtzahl $2\,p^2$, die wir so erhalten, ist noch zu halbiren, weil
hier die Substitutionen als zwei verschiedene auftreten, in denen
alle Zeichen entgegengesetzt sind; und wir erhalten also $p^2$ Substitutionen
 dieser Art (die identische eingeschlossen).

II.~\gesperrt{$(a+d)^2 - 4$ quadratischer Rest von $p$. Eine solche
Substitution lässt sich in die Normalform $S$ transformiren
mit reellem $\sigma$.}

Verstehen wir unter $\gamma$ eine primitive Wurzel des Körpers
$\mathfrak{E}_{2,p}$, so ist
\[
\gamma^{\frac{p^2-1}{2}} = -1,
\]
und wir können $r$ so bestimmen, dass
\[
\sigma = \gamma^{r(p+1)}, \quad \sigma^{\frac{p-1}{2}} = (-1)^r
\]
wird (\S.\,67). Nach (5) erhalten wir dann
\weq{9}{a + d = \gamma^{r(p+1)} + \gamma^{-r(p+1)},}
und weil nach (5) die beiden Werthe $\sigma, \sigma^{-1}$, von der Reihenfolge
abgesehen, durch $a+d$ bestimmt sind, so werden zwei Werthe
des Ausdrucks (9) nur dann einander gleich oder entgegengesetzt,
wenn die entsprechenden Werthe, mit positivem oder
%CONT>>
